\section{Minimal singularities in an arbitrary big background}
\label{sec:big-minimal-existence}

We retain the assumptions and the strict potential of \cref{sec:big-minimal-comparison}. The reference is now
\[
 V_\theta=\left(\sup\{w\in\PSH(X,\theta):w\leq0\}\right)^*.
\]
In particular $\rho\leq V_\theta\leq0$ and $\sup_XV_\theta=0$. We will solve the equations in the ordinary singularity type of $V_\theta$, with no bounded mass or uniform positive lower volume assumption. The resolution step uses the following precise version of classical regularization.

\begin{theorem}[Demailly regularization]\label{thm:resolvable-regularization-input}
Let $\alpha$ be any smooth real $(1,1)$-form and let $u\in\PSH(X,\alpha)$. There are $\delta_j\downarrow0$ and $\widetilde u_j\in\PSH(X,\alpha+\delta_j\omega)$ decreasing to $u$ such that each $\widetilde u_j$ is smooth off an analytic pole set and, after a finite sequence of blowups with smooth centers, has local expressions
\[
 \pi_j^*\widetilde u_j=\sum_\ell a_\ell\log|s_\ell|+b,
 \qquad a_\ell>0,\quad b\in C^\infty.
\]
The divisors may be taken to have simple normal crossings. The coefficients include the fixed normalization of $\ddc$.
\end{theorem}

This is the standard input \cite[Theorem 3.2]{DP04}, applied to the closed current $\ddc u$ with smooth lower bound $-\alpha$. Its lower bound is not required to be closed, and its ambient manifold is only Hermitian. Item (iv) supplies the smooth remainder after resolution; ordinary analytic type with merely a bounded remainder would not suffice here. This version need not be the particular multiplier-ideal sequence in \cref{thm:demailly-hermitian-input}. We do not assert that the additional conclusions of these two regularization statements hold for the same sequence. Applying this theorem to any strictly $\theta$-plurisubharmonic potential with a sufficiently small curvature error also justifies the choice of $\rho$ above.

\begin{lemma}\label{lem:big-resolved-equation}
Let $\psi\in\PSH(X,\theta)$ satisfy $\theta_\psi\geq a\omega$ for some $a>0$ and have the resolved smooth-remainder form in the preceding theorem. For every nonzero $0\leq f\in L^p(\mathrm{d}V)$, $p>1$, and $\lambda>0$, there exists $u\in\PSH(X,\theta)$ with $u\simeq\psi$ such that
\[
 \mu_u=\mathrm{e}^{\lambda u}f\,\mathrm{d}V.
\]
\end{lemma}
\begin{proof}
Resolve $\psi$ by $\pi:Y\to X$. The pulled-back current has the form
\[
 \pi^*\theta+\ddc\pi^*\psi=\beta+[D],
\]
where $D$ is effective and $\beta$ is smooth. Off $D$ one has $\beta\geq a\pi^*\omega$, so this inequality holds everywhere by smoothness. Thus $\beta$ is semipositive.

We check its bigness. For a blowup with smooth center and exceptional divisor $E$, choose a metric on $\mathcal O(-E)$ with curvature positive in the directions tangent to the projective fibers. The pullback of a Hermitian metric is positive in the complementary directions along $E$. Consequently its sum with a sufficiently small positive multiple of this curvature is a Hermitian metric: first check positivity along $E$, then in a neighborhood, and finally on the compact complement, where the pullback is already positive. Iterating over the blowups gives an effective real exceptional divisor $F$, with logarithmic potential $\zeta$, and a smooth Hermitian form $\kappa$ such that
\[
 \pi^*\omega+\ddc\zeta=\kappa+[F].
\]
For the iteration, pull back the previous smooth positive form, apply the one-blowup construction, and add the pulled-back previous divisor to the new positive multiple of $E$. Hence
$\beta+a\ddc\zeta\geq a\kappa$, proving bigness of $\beta$. This argument concerns Hermitian metrics and never requires a closed positive form on $Y$.

Fix a smooth positive volume form $\mathrm{d}V_Y$, and write $\pi^*\mathrm{d}V=J\,\mathrm{d}V_Y$. The function $J$ is smooth, nonnegative and bounded, and is positive off the exceptional locus. Define
\[
 H=\mathrm{e}^{\lambda\pi^*\psi}(f\circ\pi)J.
\]
It is nonzero and belongs to $L^p(\mathrm{d}V_Y)$. Indeed $\psi$ is bounded above, and change of variables gives
\[
 \int_Y|(f\circ\pi)J|^p\,\mathrm{d}V_Y
 \leq\|J\|_\infty^{p-1}\int_X|f|^p\,\mathrm{d}V.
\]
Values on the analytic exceptional sets do not affect these integrals. The semipositive big-background theorem \cref{thm:big-bounded-exponential} produces a bounded $s\in\PSH(Y,\beta)$ with
\[
 \mathrm{MA}_\beta(s)=\mathrm{e}^{\lambda s}H\,\mathrm{d}V_Y.
\]
The function $\pi^*\psi+s$ is $\pi^*\theta$-plurisubharmonic, because its current is $\beta_s+[D]\geq0$. By \cref{thm:hermitian-space-descent} it descends to $u\in\PSH(X,\theta)$. Boundedness of $s$ gives $u\simeq\psi$. Off the exceptional and polar sets, change of variables proves the required equation. Both its extended left-hand side and its smooth-volume right-hand side ignore those analytic sets. The equation therefore holds on all of $X$.
\end{proof}

\begin{theorem}\label{thm:big-minimal-exponential}
Let $\theta$ be any smooth big real $(1,1)$-form and let $0\leq f\in L^p(\mathrm{d}V)$, $p>1$, be nonzero. For every $\lambda>0$ there is a unique $u\in\PSH(X,\theta)$ such that
\[
 u\simeq V_\theta,\qquad
 \mathrm{MA}_\theta(u)=\mathrm{e}^{\lambda u}f\,\mathrm{d}V.
\]
No bounded mass or uniform positive lower volume hypothesis is needed.
\end{theorem}
\begin{proof}
Apply \cref{thm:resolvable-regularization-input} to $V_\theta$. Set $\delta_j=2^{-j}$ and take a subsequence whose curvature errors are at most $\delta_j$. Denote its potentials by $\widetilde\psi_j\downarrow V_\theta$; then
$\theta_{\widetilde\psi_j}\geq-\delta_j\omega$. Set
\[
 a_j=(1-\delta_j)\sup_X\widetilde\psi_j\geq0,
 \qquad
 \psi_j=(1-\delta_j)\widetilde\psi_j+\delta_j\rho-a_j.
\]
Normalized compactness, or monotonicity and Hartogs' lemma, gives $a_j\to0$. Resolving the singularities of $\widetilde\psi_j$ and $\rho$ together shows that $\psi_j$ has a smooth remainder after resolution. Moreover,
\begin{equation}\label{eq:big-minimal-approximation}
 \theta_{\psi_j}\geq\delta_j^2\omega,
 \qquad
 (1-\delta_j)V_\theta+\delta_j\rho-a_j
 \leq\psi_j\leq V_\theta.
\end{equation}
The upper inequality follows because $\psi_j$ is $\theta$-plurisubharmonic and at most zero. The explicit subtraction of $a_j$ keeps the normalization separate from the fixed pole barrier.

By the resolved equation lemma there are $u_j\simeq\psi_j$ satisfying
\[
 \mu_{u_j}=\mathrm{e}^{\lambda u_j}f\,\mathrm{d}V.
\]
Each $u_j$ is bounded below by $\rho$ up to a constant depending on $j$. The Gauduchon trace proposition gives a uniform upper bound for $\sup_Xu_j$.

There is also a uniform lower estimate. Translate the fixed $\rho$ down by $1+\sup_j a_j$ when applying \cref{lem:big-reference-subsolution}, and take its reference to be $\psi_j-1$. The lemma gives $\psi_j-1\leq v_j\leq\psi_j\leq0$ and $\mu_{v_j}\geq mf\,\mathrm{d}V$, with $m>0$ independent of $j$. Put $w_j=v_j+\lambda^{-1}\log m$. Then
\[
 \mu_{w_j}\geq \mathrm{e}^{\lambda w_j}f\,\mathrm{d}V,
 \qquad u_j\simeq w_j\simeq\psi_j.
\]
The exponential order in \cref{cor:big-barrier-orders} gives $u_j\geq w_j$. Thus for a uniform $C$,
\begin{equation}\label{eq:big-minimal-solution-bounds}
 \psi_j-C\leq u_j\leq V_\theta+C.
\end{equation}
After extraction, $u_j$ converges in $L^1$ and almost everywhere to a genuine $u$. The preceding bounds give $V_\theta-C\leq u\leq V_\theta+C$.

\Needspace{14\baselineskip}
We verify its equation using envelopes. Put $A_j=\sup_{k\geq j}a_k$ and, for $N\geq j$, define
\[
 p_{j,N}=P_\theta(u_j,\ldots,u_N),\qquad
 t_{j,N}=\max_{j\leq k\leq N}u_k.
\]
Both functions lie above the genuine competitor
\[
 b_j=(1-\delta_j)V_\theta+\delta_j\rho-C-A_j.
\]
In particular every envelope is genuine. All these functions are uniformly locally bounded on $\Omega$, also as $j$ varies. Contact support from \cref{thm:big-general-envelope-contact} and a disjoint partition by the first contact index give
\[
 \mu_{p_{j,N}}\leq \mathrm{e}^{\lambda p_{j,N}}f\,\mathrm{d}V.
\]
There is no factor $N$ here: on each piece the envelope equals one $u_k$, and ordered contact bounds its measure by that single measure. Locality and ordered contact for a finite maximum similarly give
\[
 \mu_{t_{j,N}}\geq \mathrm{e}^{\lambda t_{j,N}}f\,\mathrm{d}V.
\]
Let $N\to\infty$, setting $p_j=\lim_Np_{j,N}$ and $t_j=(\lim_Nt_{j,N})^*$. Local bounded monotone continuity passes both inequalities to these limits. Next $p_j$ increases to a regularized limit $p$, while $t_j$ decreases to $u$ by Hartogs' lemma. The bound by $b_j$ gives $p\simeq V_\theta$, and $p\leq u$. A second use of local bounded monotone continuity gives
\[
 \mu_p\leq \mathrm{e}^{\lambda p}f\,\mathrm{d}V,
 \qquad \mu_u\geq \mathrm{e}^{\lambda u}f\,\mathrm{d}V.
\]
These inequalities hold on $\Omega$ and extend across $Z$ because all measures involved ignore it. Thus
$\mu_p\leq \mathrm{e}^{\lambda(p-u)}\mu_u$. The exponential order gives $p\geq u$, so $p=u$ and the equation follows. Applying that order in both directions also proves uniqueness. Every passage of the Monge--Amp\`ere operator above was along a locally bounded monotone sequence on $\Omega$.
\end{proof}

\begin{theorem}\label{thm:big-minimal-normalized}
Under the assumptions of \cref{thm:big-minimal-exponential}, there are $v\in\PSH(X,\theta)$ and a uniquely determined $c>0$ such that
\[
 \sup_Xv=0,
 \qquad V_\theta-C\leq v\leq V_\theta,
 \qquad \mathrm{MA}_\theta(v)=cf\,\mathrm{d}V.
\]
Uniqueness of the normalized potential is not asserted.
\end{theorem}
\begin{proof}
Solve the positive-exponent equations with $\lambda=1/j$, and write
\[
 v_j=u_j-\sup_Xu_j,\qquad
 c_j=\mathrm{e}^{\sup_Xu_j/j},\qquad
 \mu_{v_j}=c_j \mathrm{e}^{v_j/j}f\,\mathrm{d}V.
\]
Then $v_j\simeq V_\theta$ and $v_j\leq V_\theta\leq0$. The Gauduchon trace gives $c_j\leq C_0$ uniformly.

\Needspace{12\baselineskip}
Choose $\varepsilon>0$ small and $1<q<p$ such that $\mathrm{e}^{-\varepsilon v_j}f$ is bounded in $L^q$. This follows from H\"older's inequality and uniform Skoda integrability of normalized potentials in a fixed enlarged Hermitian background. The subsolution lemma, with reference $V_\theta-1$ and the fixed barrier $\rho-1$, produces
\[
 V_\theta-1\leq\eta_j\leq V_\theta,
 \qquad
 \mu_{\eta_j}\geq m \mathrm{e}^{-\varepsilon v_j}f\,\mathrm{d}V
                 \geq mf\,\mathrm{d}V
\]
for a uniform $m>0$. Since $\mu_{v_j}\leq c_jf\,\mathrm{d}V$, scalar order gives $c_j\geq m$.

\Needspace{8\baselineskip}
Take $D=\varepsilon^{-1}\max\{0,\log(C_0/m)\}$ and put $w_j=\eta_j-D$. We have
\[
 \mathrm{e}^{\varepsilon(w_j-v_j)}\mu_{v_j}
 =c_j \mathrm{e}^{v_j/j+\varepsilon\eta_j-\varepsilon D-\varepsilon v_j}
           f\,\mathrm{d}V
 \leq m \mathrm{e}^{-\varepsilon v_j}f\,\mathrm{d}V
 \leq\mu_{w_j}.
\]
Exponential order therefore gives $v_j\geq w_j$. In particular
\begin{equation}\label{eq:big-normalized-uniform-type}
 V_\theta-1-D\leq v_j\leq V_\theta,
 \qquad m\leq c_j\leq C_0.
\end{equation}
Extract a subsequence with $v_j\to v$ in $L^1$ and almost everywhere and $c_j\to c\in[m,C_0]$. The limit is normalized, genuine and of minimal ordinary type. After relabeling, denote the retained exponents by $\lambda_j\downarrow0$; thus $\mu_{v_j}=c_j\mathrm{e}^{\lambda_jv_j}f\,\mathrm{d}V$.

One more comparison is needed to identify the measure. Define finite lower envelopes $p_{j,N}=P_\theta(v_j,\ldots,v_N)$ and finite maxima $t_{j,N}=\max_{j\leq k\leq N}v_k$. They all lie between $V_\theta-1-D$ and $V_\theta$. On $\Omega$ put
\[
 h_k=c_k \mathrm{e}^{(\lambda_k-\varepsilon)v_k},\qquad
 H_j=\sup_{k\geq j}h_k,\qquad L_j=\inf_{k\geq j}h_k.
\]
These functions are bounded on each relatively compact subset of $\Omega$, uniformly in $j$. Almost everywhere they converge, as $j\to\infty$, to $c \mathrm{e}^{-\varepsilon v}$. Contact partition and the maximum principle give
\[
 \mu_{p_{j,N}}\leq \mathrm{e}^{\varepsilon p_{j,N}}H_jf\,\mathrm{d}V,
 \qquad
 \mu_{t_{j,N}}\geq \mathrm{e}^{\varepsilon t_{j,N}}L_jf\,\mathrm{d}V.
\]
First let $N\to\infty$, obtaining $p_j$ and $t_j$, and then let $j\to\infty$. The lower envelopes increase to $p\simeq V_\theta$ with $p\leq v$, while the upper envelopes decrease to $v$. Local bounded monotone continuity applies to the left sides, and dominated convergence on relatively compact subsets of $\Omega$ applies to the densities. We obtain
\[
 \mu_p\leq c \mathrm{e}^{\varepsilon(p-v)}f\,\mathrm{d}V,
 \qquad \mu_v\geq cf\,\mathrm{d}V.
\]
Extend across $Z$ and apply exponential order to get $p\geq v$. Hence $p=v$ and $\mu_v=cf\,\mathrm{d}V$. The artificial fixed exponent $\varepsilon$ in this last argument is what allows domination to identify the two envelope limits; the unweighted inequalities alone would not do so. Scalar uniqueness is \cref{cor:big-barrier-orders}.
\end{proof}

These theorems reconstruct \cite[Theorems 4.6--4.7]{BGL25}. Their equations prove finite mass for the solutions, not a mass bound for their whole ordinary singularity type. At exponent zero only the scalar is unique; \cref{thm:degenerate-hermitian-nonuniqueness} gives counterexamples to uniqueness of the potential.
